\documentclass[11pt]{amsart}
\usepackage[T1]{fontenc}
\usepackage{lmodern,microtype,amsmath,amssymb,amsthm,mathtools,enumitem,booktabs}
\usepackage[margin=1.1in]{geometry}
\usepackage[hidelinks]{hyperref}
\usepackage{xurl}

\newtheorem{theorem}{Theorem}[section]
\newtheorem{lemma}[theorem]{Lemma}
\newtheorem{corollary}[theorem]{Corollary}
\newtheorem{proposition}[theorem]{Proposition}
\newtheorem{conjecture}[theorem]{Conjecture}
\theoremstyle{definition}
\newtheorem{definition}[theorem]{Definition}
\theoremstyle{remark}
\newtheorem{remark}[theorem]{Remark}
\numberwithin{equation}{section}

\newcommand{\R}{\mathbb{R}}
\newcommand{\E}{\mathbb{E}}

\hypersetup{pdftitle={Layered permutation copulas are closer to the uniform law than the diagonal},pdfauthor={Yitzchak Shmalo}}

\title[Layered permutations and the farthest copula]{Layered permutation copulas are closer to the uniform law than the diagonal}
\author{Yitzchak Shmalo}
\address{Einstein Institute of Mathematics, The Hebrew University of Jerusalem, Jerusalem, Israel}
\email{yitzchak.shmalo@gmail.com}
\date{7 November 2026}

\begin{document}

\begin{abstract}
Schrott conjectured that among all copulas on the unit square, the diagonal and the antidiagonal couplings are the
farthest from the uniform law in the quadratic Wasserstein distance, at squared distance $1/10$. By density of
permutation copulas the conjecture is implied by a statement about $n$ equally weighted grid points: for every
permutation $\sigma$ of $\{0,\dots,n-1\}$, the points $\bigl((i+\frac12)/n,(\sigma(i)+\frac12)/n\bigr)$ are no farther
from the uniform law than the diagonal points. We prove this, with strict inequality, for every layered permutation,
that is, every permutation obtained by cutting $0,\dots,n-1$ into runs of consecutive values and reversing some of the
runs, other than the identity and the full reversal. The proof re-cuts the optimal partition of the identity on the
union of the moved strips, writes the gain as a weighted sum of two-cell comparisons, and settles each comparison by a
dispersive-order inequality between the laws of $x_1+x_2$ and $x_1-x_2$ on a diagonal band. The same re-cut, done
in two steps inside a window of three or four strips, gives the strict inequality for every permutation that splits
into blocks of consecutive values of length at most four (each block reversed or permuted arbitrarily), once $n\ge9$;
for blocks of length four the positivity of six explicit functionals is proved with interval arithmetic.
\end{abstract}

\maketitle

\section{The problem}

Let $\lambda$ be Lebesgue measure on $Q=[0,1]^2$. A copula is a probability measure on $Q$ with uniform marginals.
Schrott \cite{Schrott} proved that among all couplings of $n$ standard Gaussians, the ones farthest from the standard
Gaussian on $\R^n$ in $W_2$ are the signed diagonals, and stated the uniform analogue as a conjecture.

\begin{conjecture}[{\cite[Conjecture~1.1]{Schrott}}]\label{conj}
For every copula $\pi$ on $Q$, $W_2^2(\pi,\lambda)\le 1/10$, with equality for the diagonal coupling.
\end{conjecture}

For a probability measure $\mu$ on $Q$ write $V(\mu)=\max\E[X\cdot Z]$ over couplings of $X\sim\lambda$ and $Z\sim\mu$.
Since $W_2^2(\mu,\lambda)=\E|X|^2+\E|Z|^2-2V(\mu)$ and $\E|Z|^2=2/3$ for every copula, Conjecture~\ref{conj} says that
$V(\pi)\ge V(\text{diagonal})$ for every copula $\pi$.

Fix $n\ge 2$, put $a_i=(i+\frac12)/n$ for $0\le i\le n-1$, and for a permutation $\sigma$ of $\{0,\dots,n-1\}$ let
$\mu_\sigma=\frac1n\sum_i\delta_{p_i}$ with $p_i=(a_i,a_{\sigma(i)})$, and $A_n(\sigma)=V(\mu_\sigma)$. Both marginals of
$\mu_\sigma$ are the uniform law on the grid $\{a_i\}$, so $\E|Z|^2$ does not depend on $\sigma$, and
$A_n(\sigma)\ge A_n(\mathrm{id})$ says that $\mu_\sigma$ is no farther from $\lambda$ than the diagonal points. Consider
\begin{equation}\label{An}
A_n(\sigma)\ge A_n(\mathrm{id})\qquad\text{for every permutation }\sigma\text{ of }\{0,\dots,n-1\}.
\tag{$\mathrm{A}_n$}
\end{equation}
Every copula $\pi$ is a weak limit of measures $\mu_{\sigma_n}$. Indeed, let $Z^1,\dots,Z^n$ be independent with law
$\pi$, almost surely with distinct coordinates, and let $\sigma_n$ send the rank of $Z^j_1$ to the rank of $Z^j_2$
(ranks counted from $0$). Replacing each coordinate by $(\text{rank}+\frac12)/n$ moves it by at most
$\sup_t|F_n(t)-t|+\frac1{2n}$, where $F_n$ is the empirical distribution function of that coordinate, which tends to $0$
almost surely by the Glivenko--Cantelli theorem; and the empirical measure of the $Z^j$ tends weakly to $\pi$. On $Q$
weak convergence is $W_2$ convergence, and $|V(\mu)-V(\mu')|\le(\E|X|^2)^{1/2}W_2(\mu,\mu')$. Hence, if \eqref{An} holds
for all $n$, then $V(\pi)=\lim_n A_n(\sigma_n)\ge\lim_n A_n(\mathrm{id})=V(\text{diagonal})=37/60$ for every copula
$\pi$, the last limit by the same argument applied to the diagonal. So \eqref{An} for all $n$ implies
Conjecture~\ref{conj}.

The map $A_n$ is invariant under the eight symmetries of the square acting on permutation matrices, so
$A_n(\mathrm{rev})=A_n(\mathrm{id})$ for the full reversal $\mathrm{rev}(i)=n-1-i$. Numerically \eqref{An} holds
for every permutation when $n\le 8$, with equality only at $\mathrm{id}$ and $\mathrm{rev}$ (Section~\ref{sec:numerics}).

\begin{definition}
A permutation $\sigma$ is \emph{layered} if there are integers $0=k_0<k_1<\dots<k_m=n$ such that, for each $\ell$,
$\sigma$ maps $\{k_\ell,\dots,k_{\ell+1}-1\}$ onto itself either identically or by reversal.
\end{definition}

\begin{theorem}\label{thm:main}
Let $n\ge 2$ and let $\sigma$ be a layered permutation of $\{0,\dots,n-1\}$ other than $\mathrm{id}$ and $\mathrm{rev}$.
Then $A_n(\sigma)>A_n(\mathrm{id})$. The same holds for the images of $\sigma$ under the symmetries of the square.
\end{theorem}

For $n=2$ every permutation is $\mathrm{id}$ or $\mathrm{rev}$, so the theorem is about $n\ge 3$. Section~\ref{sec:windows}
extends Theorem~\ref{thm:main} to permutations built from blocks of length at most four (Theorem~\ref{thm:windows}).

\section{The identity and the re-cut}\label{sec:recut}

Write $s=x_1+x_2$ and $v=x_1-x_2$ for $x\in Q$. Under $\lambda$, $S=X_1+X_2$ has density
$L(s)=\min(s,2-s)$ on $[0,2]$, and given $S=s$ the variable $V=X_1-X_2$ is uniform on $[-L(s),L(s)]$. Let
$0=q_0<q_1<\dots<q_n=2$ be the quantiles of $S$ at levels $j/n$ and $T_j=\{x\in Q: q_j<s<q_{j+1}\}$.

\begin{lemma}\label{lem:id}
$A_n(\mathrm{id})=\sum_j a_j\int_{T_j}s\,dx$.
\end{lemma}

\begin{proof}
A coupling of $\lambda$ with $\mu_\sigma$ gives each point $p_i$ a sub-probability $\gamma_i\le\lambda$ of mass $1/n$, and
$\E[X\cdot Z]=\sum_i\int x\cdot p_i\,d\gamma_i$. For $\sigma=\mathrm{id}$, $x\cdot p_j=a_js$, and with $a_0<\dots<a_{n-1}$
the sum $\sum_j a_j\int s\,d\gamma_j$ under $\sum_j\gamma_j=\lambda$, $\gamma_j(Q)=1/n$, is largest when $\gamma_j$ is
$\lambda$ on $T_j$ (the rearrangement inequality applied to the law of $S$).
\end{proof}

Any partition of $Q$ into regions $R_0,\dots,R_{n-1}$ of area $1/n$ gives the coupling $\gamma_i=\lambda|_{R_i}$, so
$A_n(\sigma)\ge\sum_i\int_{R_i}x\cdot p_i\,dx$.

Consider a block of $r\ge 2$ consecutive values $k,\dots,k+r-1$ with $k+r\le n$, and let $\sigma$ reverse it and fix every
other value. Its moved points are $p_{k+i}=(a_{k+i},a_{k+r-1-i})$, $0\le i\le r-1$; they share
$p_1+p_2=2\bar a$ with $\bar a=a_k+(r-1)/(2n)$, and $p_1-p_2=w_i:=(2i+1-r)/n$, increasing in $i$. Keep $T_j$ for every
fixed point, and cut $U=T_k\cup\dots\cup T_{k+r-1}=\{q_k<s<q_{k+r}\}$ into the regions $R_0,\dots,R_{r-1}$ of area $1/n$
cut out by consecutive quantiles of $v$ on $U$, from the smallest values of $v$ up; give $R_i$ to $p_{k+i}$. Since
$x\cdot p=\frac12\bigl(s(p_1+p_2)+v(p_1-p_2)\bigr)$, Lemma~\ref{lem:id} gives $A_n(\sigma)-A_n(\mathrm{id})\ge D$ with
\[
D=\sum_{i=0}^{r-1}\int_{R_i}\Bigl(\bar a s+\tfrac12 w_iv\Bigr)dx-\sum_{j=k}^{k+r-1}a_j\int_{T_j}s\,dx
=\frac1{n^2}\Bigl(\sum_{i=0}^{r-1}(i-\tfrac{r-1}2)\,\nu_i-\sum_{i=0}^{r-1}(i-\tfrac{r-1}2)\,m_i\Bigr),
\]
where $m_i=n\int_{T_{k+i}}s\,dx$ is the mean of $S$ on the strip $T_{k+i}$ and $\nu_i=n\int_{R_i}v\,dx$ the mean of $V$
on $R_i$.

\begin{lemma}\label{lem:split}
For real numbers $c_0,\dots,c_{r-1}$,
$\sum_i(i-\frac{r-1}2)c_i=\sum_{t=1}^{r-1}\frac{t(r-t)}r\,\Delta_t(c)$ with
$\Delta_t(c)=\frac1{r-t}\sum_{i\ge t}c_i-\frac1t\sum_{i<t}c_i$.
\end{lemma}

\begin{proof}
With $\bar c$ the mean, $\sum_i i(c_i-\bar c)=\sum_{t\ge1}\sum_{i\ge t}(c_i-\bar c)$, and
$\sum_{i\ge t}(c_i-\bar c)=(r-t)\bigl(\frac1{r-t}\sum_{i\ge t}c_i-\bar c\bigr)=\frac{t(r-t)}r\Delta_t(c)$.
\end{proof}

So $n^2D=\sum_{t=1}^{r-1}\frac{t(r-t)}r\bigl(\Delta_t(\nu)-\Delta_t(m)\bigr)$. Both $\nu$ and $m$ are the means of a
variable over its own quantile cells of mass $1/n$ on $U$: $\Delta_t(\nu)$ is the gap between the mean of $V$ above and
below its quantile at mass $t/n$ on $U$, and $\Delta_t(m)$ the same gap for $S$.

\section{A dispersive comparison on diagonal bands}

Fix $0\le\alpha<\beta\le2$ with $(\alpha,\beta)\ne(0,2)$ and let $U=\{\alpha<s<\beta\}$, of area $\rho$. Let $\Phi_S$ and
$\Phi_V$ be the distribution functions of $S$ and $V$ under $\lambda|_U$ (total mass $\rho$), $Q_S,Q_V$ their inverses on
$(0,\rho)$, and $\phi_S=L\circ Q_S$, $\phi_V=g\circ Q_V$ the densities read at mass level $\mu$, where
$g(y)=\frac12\bigl|(\alpha,\beta)\cap(|y|,2-|y|)\bigr|$ is the density of $V$ under $\lambda|_U$. Write $\mu'=\rho-\mu$.

\begin{lemma}\label{lem:disp}
For $0<\mu<\rho$, $\phi_V(\mu)\le\sqrt{2\min(\mu,\mu')}\le\phi_S(\mu)$. If $\alpha>0$ there is $c>0$ with
$\phi_V(\mu)<\phi_S(\mu)$ for $0<\mu<c$, and if $\beta<2$ there is $c>0$ with $\phi_V(\mu)<\phi_S(\mu)$ for $\rho-c<\mu<\rho$.
\end{lemma}

\begin{proof}
Let $s=Q_S(\mu)$. If $s\le1$, then $\mu=\int_\alpha^sL\le\int_0^s t\,dt=s^2/2$, so $\phi_S(\mu)=s\ge\sqrt{2\mu}$, strictly
if $\alpha>0$. If $s\ge1$, then $\mu'=\int_s^\beta L\le(2-s)^2/2$ and $\phi_S(\mu)=2-s\ge\sqrt{2\mu'}$, strictly if
$\beta<2$. In both cases $\phi_S(\mu)\ge\sqrt{2\min(\mu,\mu')}$.

For the left inequality, $g$ is even. For $y\in[-1,0]$ the interval $(-y,2+y)$ grows with $y$ and each of its ends
moves at unit speed, so $g$ is nondecreasing on $[-1,0]$, $g(-1)=0$, and $g(y)-g(z)\le y-z$ for $-1\le z\le y\le0$. Hence
for $y=Q_V(\mu)\le0$,
\[
\mu=\int_{-1}^yg(z)\,dz\ge\int_{-1}^y\bigl(g(y)-(y-z)\bigr)_+dz\ge\tfrac12g(y)^2,
\]
so $\phi_V(\mu)\le\sqrt{2\mu}$, and $\mu\le\mu'$ by the symmetry of $g$. The case $y\ge0$ is the mirror image.

For the strict inequality let $\alpha>0$. If $\alpha<1$, put $c=\min(\Phi_S(1),\rho)>0$: for $\mu<c$ we have $Q_S(\mu)<1$
and $\phi_S(\mu)>\sqrt{2\mu}\ge\phi_V(\mu)$ by the first paragraph. If $\alpha\ge1$, put $c=\rho/2$: then $Q_S(\mu)\ge1$, and
for $\mu<c$ we have $\mu<\mu'$, so $\phi_S(\mu)\ge\sqrt{2\mu'}>\sqrt{2\mu}\ge\phi_V(\mu)$. The case $\beta<2$ follows by
applying the map $(x_1,x_2)\mapsto(1-x_2,1-x_1)$, which preserves $\lambda$, fixes $v$, sends $s$ to $2-s$ and the band
$(\alpha,\beta)$ to $(2-\beta,2-\alpha)$, and reverses the mass levels of $S$.
\end{proof}

\begin{proposition}\label{prop:block}
Let $n\ge3$, $2\le r\le n-1$, $0\le k\le n-r$. Then $D>0$ for the reversal of the block $k,\dots,k+r-1$.
\end{proposition}

\begin{proof}
Take $\alpha=q_k$, $\beta=q_{k+r}$; since $r<n$, $(\alpha,\beta)\ne(0,2)$. Since $Q_S'=1/\phi_S$ and
$Q_V'=1/\phi_V$, Lemma~\ref{lem:disp} shows that $h=Q_V-Q_S$ is nondecreasing on $(0,\rho)$. For $\tau=t/n$,
\[
\Delta_t(\nu)-\Delta_t(m)=\frac1{\tau(\rho-\tau)}\int_0^\tau\!\!\int_\tau^\rho\bigl(h(\mu)-h(\mu_0)\bigr)\,d\mu\,d\mu_0\ge0.
\]
If $\alpha>0$, Lemma~\ref{lem:disp} gives $c>0$ with $h$ strictly increasing on $(0,c)$, so $h(\mu)>h(\mu_0)$ whenever
$\mu_0<\min(\tau,c)/2$ and $\mu>\tau$; the integral is then positive for every $t$. If $\alpha=0$ then $\beta<2$, $h$ is
strictly increasing on some $(\rho-c,\rho)$, and $h(\mu)>h(\mu_0)$ whenever $\mu_0<\tau$ and
$\mu>\rho-\min(\rho-\tau,c)/2$. By Lemma~\ref{lem:split}, $n^2D>0$.
\end{proof}

\begin{proof}[Proof of Theorem~\ref{thm:main}]
Let the reversed runs of $\sigma$ be $B_1,\dots,B_m$, of sizes $r_\ell\ge2$. If some $r_\ell=n$ then $\sigma=\mathrm{rev}$;
otherwise each $r_\ell\le n-1$, and $m\ge1$ since $\sigma\ne\mathrm{id}$. The unions $U_\ell$ of the strips of the runs are
disjoint, so performing the re-cut of Section~\ref{sec:recut} inside every $U_\ell$ and keeping the other strips is a partition of $Q$
into regions of area $1/n$ matched to the points of $\mu_\sigma$. Its value is $A_n(\mathrm{id})+\sum_\ell D_\ell$, and
each $D_\ell>0$ by Proposition~\ref{prop:block}. The symmetries of the square preserve $\lambda$ and the grid, so $A_n$
is invariant under them.
\end{proof}

\begin{remark}
For $r=2$ at the corner ($k=0$ or $k=n-2$, $n\ge4$) the bound is exactly
$n^2D=\frac{2\sqrt2-2}{3}\,n^{-1/2}\approx0.276\,n^{-1/2}$, while the measured gap $n^2(A_n(\sigma)-A_n(\mathrm{id}))$ is
about $0.344\,n^{-1/2}$ ($n=8,12,16$). The corner transposition is the cheapest move, so the margin in \eqref{An} near
the identity is of order $n^{-5/2}$.
\end{remark}

\section{Windows of three and four values}\label{sec:windows}

Let $r\ge2$, let $\tau$ be a permutation of $\{0,\dots,r-1\}$, $0\le k\le n-r$, and let $\sigma(k+i)=k+\tau(i)$ for
$0\le i\le r-1$, $\sigma$ the identity elsewhere. Put $u'_i=i+\tau(i)-(r-1)$ and $v'_i=i-\tau(i)$. The moved points
satisfy $p_1+p_2=2\bar a+u'_i/n$ and $p_1-p_2=v'_i/n$ for $p=p_{k+i}$, with $\bar a$ and $U=\{q_k<s<q_{k+r}\}$ as in
Section~\ref{sec:recut}. For a region $C$ of positive area write $\E_C$ for the average over $C$ with respect to $\lambda$.

\begin{lemma}\label{lem:window}
Let $C_0,\dots,C_{r-1}$ be a partition of $U$ into regions of area $1/n$, and
\[
F=\sum_{i=0}^{r-1}\bigl(u'_i\,\E_{C_i}[S]+v'_i\,\E_{C_i}[V]\bigr)-\sum_{j=0}^{r-1}(2j-r+1)\,\E_{T_{k+j}}[S].
\]
Then $A_n(\sigma)-A_n(\mathrm{id})\ge F/(2n^2)$. If $\sigma$ acts in this way on several disjoint blocks, re-cutting each
block's strips separately gives $A_n(\sigma)-A_n(\mathrm{id})\ge\sum_\ell F_\ell/(2n^2)$.
\end{lemma}

\begin{proof}
Give $C_i$ to $p_{k+i}$ and keep $T_j$ for the fixed points. With $x\cdot p=\frac12(s(p_1+p_2)+v(p_1-p_2))$ and
$\int_C g\,dx=\frac1n\E_C[g]$,
$\sum_i\int_{C_i}x\cdot p_{k+i}\,dx=\bar a\int_Us\,dx+\frac1{2n^2}\sum_i(u'_i\E_{C_i}[S]+v'_i\E_{C_i}[V])$, while by
Lemma~\ref{lem:id} the identity collects
$\sum_ja_{k+j}\int_{T_{k+j}}s\,dx=\bar a\int_Us\,dx+\frac1{2n^2}\sum_j(2j-r+1)\E_{T_{k+j}}[S]$.
Disjoint blocks use disjoint unions of strips, so the gains add.
\end{proof}

For the reversal of the block, $u'_i=0$ and $v'_i=2i-r+1$, and $F/(2n^2)$ is the quantity $D$ of
Section~\ref{sec:recut}. The strips $T_{k+j}$ are the cells of $U$ between consecutive quantiles of $S$. We use two
partitions. \emph{$v$ then $s$}: group the indices $i$ by the value of $v'_i$; cut $U$ by $V$ at its quantiles of mass
$1/n$ times the cumulative group sizes, the lowest values of $V$ going to the group with the smallest $v'$; cut each
group's region by $S$ into cells of area $1/n$ and give them to the group's indices in increasing order of $u'_i$.
\emph{$s$ then $v$} is the same with the roles of $(S,u')$ and $(V,v')$ exchanged. The band $U$ is symmetric under
$v\mapsto-v$, so the median of $V$ on $U$ is $0$.

The symmetries $(x_1,x_2)\mapsto(x_2,x_1)$ and $(x_1,x_2)\mapsto(1-x_2,1-x_1)$ preserve $\lambda$, map strips to strips
and windows to windows, and replace $\tau$ by $\tau^{-1}$ and by $i\mapsto r-1-\tau^{-1}(r-1-i)$; their composite gives
$i\mapsto r-1-\tau(r-1-i)$. The second sends the band $\{\alpha<s<\beta\}$ to $\{2-\beta<s<2-\alpha\}$. The image of a partition is a partition with the
same value of $F$. So it suffices to treat one block in each orbit, on every band.

\begin{proposition}\label{prop:r3}
Let $n\ge4$ and $\tau=(1,2,0)$ or $(2,0,1)$. Then the $v$ then $s$ partition gives $F>0$ on every band $U$ of three
consecutive strips.
\end{proposition}

\begin{proof}
The two 3-cycles form one orbit; take $\tau=(1,2,0)$, with $(u'_i,v'_i)=(-1,-1),(1,-1),(0,2)$. Let $c$ be the upper
tercile of $V$ on $U$ (so $c\ge0$), $C_2=U\cap\{v>c\}$, and split the rest $U'=U\setminus C_2$ at the median of $S$ on
$U'$ into $C_0$ (lower) and $C_1$ (upper). Since $\E_U[V]=0$, $\E_{C_0}[V]+\E_{C_1}[V]=2\E_{U'}[V]=-\E_{C_2}[V]$, and
\[
F=3\,\E_{C_2}[V]+\bigl(\E_{C_1}[S]-\E_{C_0}[S]\bigr)-2\bigl(\E_{T_{k+2}}[S]-\E_{T_k}[S]\bigr)
=\tfrac32(G_V-G_S)+\bigl(\Delta_{U'}(S)-\tfrac12G_S\bigr),
\]
where $G_X$ is the mean of $X$ on the upper third of $U$ (by mass) minus its mean on the lower third, and
$\Delta_{U'}(S)$ the gap between the means of $S$ above and below its median on $U'$. With $h=Q_V-Q_S$ nondecreasing
(Lemma~\ref{lem:disp}), $G_V-G_S=\frac3\rho\bigl(\int_{2\rho/3}^\rho h-\int_0^{\rho/3}h\bigr)\ge0$. For the second
bracket, the section of $U'$ over $s$ has length $L(s)+\min(L(s),c)\ge L(s)$, half the section of $U$. For any law,
the gap between the means above and below the median is $2\min_t\E|X-t|$; the mass of $U'$ is $2/3$ that of $U$, so
$\Delta_{U'}(S)\ge\frac34\Delta_U(S)$. The upper half of $U$ is its upper third together with the upper half of its
middle third, in proportions $2:1$, and likewise below, so $\Delta_U(S)>\frac23G_S$. Hence
$\Delta_{U'}(S)-\frac12G_S>0$ and $F>0$.
\end{proof}

For $r=4$ there are thirteen indecomposable permutations (no proper initial segment $\{0,\dots,j-1\}$ is mapped onto
itself); one is the reversal and the other twelve fall into six orbits,
\[
\{1302,2031\},\ \{1320,2130,3021,3102\},\ \{3120\},\ \{2301\},\ \{2310,3201\},\ \{1230,3012\},
\]
writing $\tau$ as the word $\tau(0)\tau(1)\tau(2)\tau(3)$. On a band $U$ let $S_1,\dots,S_4$ be the means of $S$ on the
four strips, $d_{41}=S_4-S_1$, $d_{32}=S_3-S_2$, $x_m$ the median of $|V|$ on $U$, $A$ and $B$ the means of $|V|$ above and
below $x_m$, $T_{\mathrm{out}}$ and $T_{\mathrm{in}}$ the means of $S$ on $\{|v|>x_m\}$ and $\{|v|<x_m\}$, and
$R_1,R_3$ the means of $S$ on the lowest and highest thirds (by mass, in $s$) of $\{v<x_m\}$. By the symmetry
$v\mapsto-v$ the quartiles of $V$ on $U$ are $-x_m,0,x_m$.

\begin{lemma}\label{lem:r4forms}
The $v$ then $s$ partition gives
\begin{align*}
F(1302)&=4A+2B-3d_{41}-d_{32}, & F(1320)&=5A+B+(T_{\mathrm{out}}-T_{\mathrm{in}})-3d_{41}-d_{32},\\
F(2301)&=4A+4B-2d_{41}, & F(2310)&=5A+3B-\tfrac52d_{41}-\tfrac12d_{32},\\
F(1230)&=4A+2(R_3-R_1)-3d_{41}-d_{32}, & F(3120)&=6A+(T_{\mathrm{hi}}-T_{\mathrm{lo}})-3d_{41}-d_{32},
\end{align*}
where $T_{\mathrm{hi}},T_{\mathrm{lo}}$ are the means of $S$ above and below its median on $\{|v|<x_m\}$.
\end{lemma}

\begin{proof}
For $1302$ the points are $(u',v')=(-2,-1),(1,-2),(-1,2),(2,1)$; the $v'$ are distinct, so the cells are the quartile
cells of $V$, $\{v<-x_m\},\{-x_m<v<0\},\{0<v<x_m\},\{v>x_m\}$, carrying $v'=-2,-1,1,2$ and $u'=1,-2,2,-1$. The $V$-means are
$-A,-B,B,A$, which gives $4A+2B$; the $S$-means are $T_{\mathrm{out}},T_{\mathrm{in}},T_{\mathrm{in}},T_{\mathrm{out}}$,
and the weights $1-1$ and $-2+2$ cancel them. The strips contribute $-(3S_4+S_3-S_2-3S_1)=-3d_{41}-d_{32}$. The other five are the same
computation. For 1320, $(u',v')=(-2,-1),(1,-2),(1,0),(0,3)$, the cells are again the $V$-quartiles. For 2301 the points
are $(\pm1,\pm2)$: the halves $\{v<0\}$, $\{v>0\}$ are split at the $S$-median, which is that of $U$, and the $V$-means
add to $2(A+B)$ on each side. For 2310, $(u',v')=(-1,-2),(1,-2),(0,1),(0,3)$, the half $\{v<0\}$ is split by $S$ and the
upper half into $\{0<v<x_m\}$ and $\{v>x_m\}$. For 1230, $(u',v')=(-2,-1),(0,-1),(2,-1),(0,3)$, the top quarter is
$\{v>x_m\}$ with $V$-mean $A$, the rest $\{v<x_m\}$ has $V$-mean $-A/3$ and is cut into $S$-thirds. For 3120,
$(u',v')=(0,-3),(-1,0),(1,0),(0,3)$, the outer quarters go to $v'=\pm3$ and $\{|v|<x_m\}$ is split at its $S$-median.
\end{proof}

\begin{proposition}\label{prop:r4}
Let $\tau$ be one of the twelve indecomposable non-layered permutations of $\{0,1,2,3\}$. On every band $U$ of mass
$\rho\le4/9$, the $v$ then $s$ partition gives $F>0$.
\end{proposition}

\begin{proof}
By the orbits it suffices to treat the six representatives. Always $A\ge x_m\ge B\ge0$, $T_{\mathrm{hi}}>T_{\mathrm{lo}}$,
$R_3>R_1$ and $d_{41}>d_{32}>0$. Lemma~\ref{lem:r4forms} gives
\begin{align*}
F(3120)-F(1302)&=2(A-B)+(T_{\mathrm{hi}}-T_{\mathrm{lo}}),\qquad F(2301)-F(1302)=2B+d_{41}+d_{32},\\
F(2310)-F(1302)&=A+B+\tfrac12(d_{41}+d_{32}),\qquad F(1230)\ge F_0:=4A-3d_{41}-d_{32},
\end{align*}
all differences positive, and $F(1302)=F_0+2B\ge F_0$. So it remains to show $F_0>0$ and $F(1320)>0$ for
$\rho\le4/9$. Write $U=\{\alpha<s<\beta\}$ and $w=\beta-\alpha$.

\emph{Thin bands across the diagonal}, $\alpha<1<\beta$, $w<1/3$. On $U$, $L\ge1-w$, so $\rho\ge w(1-w)$. The
distribution function of $|V|$ on $U$ is $H(x)=\int_\alpha^\beta\min(L(s),x)\,ds\le wx$, and $H(x_m)=\rho/2$, so
$x_m\ge(1-w)/2$ and $A\ge(1-w)/2$. All $S_j$, $T$ and $R$ lie in $(\alpha,\beta)$, so each difference of them is at
most $w$ in absolute value. Hence $F_0\ge2(1-w)-4w>0$ and $F(1320)\ge\frac52(1-w)-5w>0$.

\emph{Bands on one side of the diagonal}, $\beta\le1$ (the case $\alpha\ge1$ follows by $s\mapsto2-s$, which leaves $A$,
$B$, $d_{41}$, $d_{32}$ unchanged and reverses the sign of $T_{\mathrm{out}}-T_{\mathrm{in}}$). Every quantity in the
functionals is a length, and the band $\{\alpha<s<\beta\}$ is the band $\{\alpha/\beta<s<1\}$ scaled by $\beta$, so
each functional equals $\beta G(\alpha/\beta)$ for an explicit algebraic function $G$ on $[0,1)$: the quantiles of $S$
and of $|V|$ solve quadratic equations, and the means are ratios of polynomial integrals. We enclosed $G$ on
subintervals of $[0,1]$ in interval arithmetic with outward rounding (\texttt{mpmath.iv}), bisecting until the lower
end of the enclosure is positive: $F_0$ in 24 boxes (smallest lower bound $0.035$), $F(1320)$ in 12 boxes on
$\beta\le1$ ($0.124$) and 17 boxes on $\alpha\ge1$ ($0.172$). Since this covers every
$t=\alpha/\beta\in[0,1]$, it holds for bands of any mass.

\emph{Wider bands across the diagonal}, $\alpha<1<\beta$, $\rho\le4/9$, $w\ge w_0$. In the coordinates $(\alpha,2-\beta)\in[0,1]^2$,
$\rho=(2-\alpha^2-(2-\beta)^2)/2$, and $\rho F$ is a combination of the integrals $K(x)=\int_0^xsL(s)\,ds$ at
$\alpha$, $\beta$ and the quartiles of $S$, of $N(x)=\int_0^xy\,dH(y)$ at $x_m$ and $1$, and, for $1320$, of
$P(x_m)=\int_\alpha^\beta s\,(L(s)-x_m)_+\,ds$, which gives $T_{\mathrm{out}}$. The quartiles of $S$ increase in $\alpha$
and $\beta$, $K$ and $N$ increase and $P$ decreases in $x$, and $x_m$ is enclosed by bisection on the sign of $H(x)-\rho/2$ over a whole box; so each
box of parameters gets an enclosure of $\rho F$. Bisecting boxes until the enclosure is positive: $F_0$ with $w_0=0.05$ in 2,774 boxes, $F(1320)$ with $w_0=0.25$ in 1,042 boxes, no failures.
Together with the thin case this covers every band of mass at most $4/9$.
\end{proof}

\begin{theorem}\label{thm:windows}
Let $n\ge9$, and let $\sigma\ne\mathrm{id},\mathrm{rev}$ map each interval $\{k_\ell,\dots,k_{\ell+1}-1\}$ of a cut
$0=k_0<\dots<k_m=n$ onto itself, acting on each interval either by reversal or by a permutation of an interval of
length at most four. Then $A_n(\sigma)>A_n(\mathrm{id})$, and the same holds for the images of $\sigma$ under the
symmetries of the square. If every interval of length four is reversed or fixed, $n\ge4$ suffices.
\end{theorem}

\begin{proof}
Refine the cut so that $\sigma$ is indecomposable on each interval. An indecomposable piece of length at most four is
a fixed point, a reversal, one of the two 3-cycles, or one of the twelve blocks of Proposition~\ref{prop:r4}; a reversed
piece of length $n$ would make $\sigma=\mathrm{rev}$. By Lemma~\ref{lem:window} the gain is at least the sum over the
moved pieces of $F_\ell/(2n^2)$, and each $F_\ell$ is positive: reversals by Proposition~\ref{prop:block}, 3-cycles by
Proposition~\ref{prop:r3}, the twelve blocks by Proposition~\ref{prop:r4}, since four strips have mass $4/n\le4/9$.
\end{proof}

For $n\le8$ every permutation was checked numerically (Section~\ref{sec:numerics}). The blocks $1302$ and $1320$ need the
restriction on the mass: their $v$ then $s$ value is negative on nearly full bands, $\{0<s<1.8\}$ for $1302$ and $\{0.2<s<2\}$ for $1320$.

\section{Computations}\label{sec:numerics}

All computations use exact moments of $S$ and one-dimensional quadrature of the law of $V$ on a band.
\begin{enumerate}[leftmargin=*]
\item Check of \eqref{An} over every permutation for $n\le8$ (one per class under the symmetries of the square;
5,282 classes at $n=8$), with $A_n$ computed by a smoothed semi-discrete dual on a $40\times40$ grid of cells.
No permutation falls below the identity; the smallest positive gap is the
corner adjacent transposition ($n^2(A_n(\sigma)-A_n(\mathrm{id}))=0.1215$ at $n=8$), and the first non-layered
permutations are corner 3-cycles ($0.3944$ at $n=8$).
\item The bound $D$ for every block reversal, $3\le n\le48$, $2\le r\le n-1$, every $k$: all positive, with minimum at the
corner adjacent transposition, $n^{5/2}D=0.2761$, matching the closed form above. Quadrature error below $10^{-8}$.
\item The identity of Lemma~\ref{lem:split} applied to $D$, checked against the direct evaluation for $n=5,8,13$ and
every $(k,r)$: difference below $1.1\times10^{-8}$; every term $\Delta_t(\nu)-\Delta_t(m)$ is positive.
\item The closed forms of Lemma~\ref{lem:r4forms} against a direct evaluation of the $v$ then $s$ cut on polygons, on seven
bands: agreement to $10^{-7}$; the values agree across each orbit to $5\times10^{-15}$.
\item The interval enclosures of Proposition~\ref{prop:r4}, in \texttt{mpmath.iv}; the scripts are in the repository. As an
independent check, $F(1302)$ was also proved positive directly: 17 boxes on one-sided bands (smallest lower bound
$0.129$) and 1,550 boxes on crossing bands of mass at most $4/9$.
\end{enumerate}

\section{What remains}

Theorems~\ref{thm:main} and~\ref{thm:windows} cover the permutations that split into layers and blocks of at most four
values. Inside a window of five strips a directional cut still gives positive bounds for every indecomposable block that
was tested ($n=8$), but no closed form of such a cut is known for blocks of five or more values. On the whole square a cut along one direction is not
enough: cutting $Q$ into quantile cells of $s\cos\theta+v\sin\theta$ and matching the cells optimally never beats the
strips for any $\theta$ when $n=4$. A proof of \eqref{An} along these lines would need cuts whose direction varies across
the square. For $n\le8$ the inequality rests on the computation of Section~\ref{sec:numerics}.

\section*{Acknowledgements and funding}
The research presented in this paper was supported by the European Research
Council (ERC) under the European Union's Horizon Europe research and
innovation programme (grant agreement No.~101041711), by the Simons
Foundation as part of the Collaboration on the Mathematical and Scientific
Foundations of Deep Learning, by Heights Labs, by Convex Nexus Capital, by
the Israel Science Foundation (grant number 2258/19), and by the Israel
Science Foundation (ISF Grant 4101/25). Computational resources were
provided in part by Google Cloud through its education credits program
(credit 530550103).

\section*{Declaration of competing interest}
The author holds equity in Heights Labs and in Convex Nexus Capital and
has a management role at Convex Nexus Capital. These entities are included
in the funding declaration above.

\section*{Declaration of generative AI and AI-assisted technologies}
The paper comes out of Math Brain, an ongoing project on the use of AI systems for research mathematics. In November 2026
Claude Opus~5.5 (Anthropic), working in Claude Code, found the reduction, the re-cut construction and the proofs, wrote and
ran the computations of Section~\ref{sec:numerics}, and wrote the paper. The author directed the work and is responsible
for the content of the paper and for decisions concerning its dissemination.

\section*{Data availability}
No data were used. The code for the computations of Sections~\ref{sec:windows} and~\ref{sec:numerics}, including the
interval-arithmetic proofs, with its output, is available at
\url{https://github.com/yspennstate/wasserstein-layered-permutation-copulas}.

\begin{thebibliography}{9}
\bibitem{Schrott} S.~Schrott, \emph{Couplings farthest from the independent Gaussian}, arXiv:2609.10467 (2026).
\end{thebibliography}

\end{document}
